
\newgeometry{left=3cm,right=3cm,top=2.5cm,bottom=2cm,includefoot}
\fontsize{12pt}{14.4pt}\selectfont
\begin{abstract}

Орчин үеийн боловсролын орчинд мэдээллийн хэмжээ эрчимтэй өсөж, суралцагчид богино хугацаанд их хэмжээний агуулга боловсруулах шаардлагатай болж байна. Ялангуяа их, дээд сургуулийн түвшинд оюутнууд нэг хичээлийн хүрээнд лекцийн тэмдэглэл, сурах бичиг, өгүүлэл, гарын авлага, даалгаврын материал зэрэг олон төрлийн эх сурвалжтай зэрэг ажиллах хэрэгцээтэй байдаг. Гэвч эдгээр материалыг үр дүнтэй уншиж ойлгох, гол санааг ялгах, хэрэгтэй хэсгийг хурдан олох, тухайн агуулгыг өөрийн мэдлэгийн түвшинтэй уялдуулан суралцах нь бодит байдал дээр амаргүй байдаг. Ийм нөхцөлд оюутан зөвхөн мэдээлэлтэй байх бус, тухайн мэдээллийг ойлгож, холбож, хэрэглэж сурах дэмжлэг илүү их шаардлагатай болж байна.

Монголын боловсролын салбарт ч энэ асуудал тодорхой ажиглагддаг. Оюутнуудын хувьд хичээлийн материал их боловч түүнийг тайлбарлаж өгөх, тухайн сэдвийн чухал хэсгийг ялгаж өгөх, ойлгоогүй хэсэг дээр нь дахин чиглүүлэх ухаалаг дэмжлэг дутмаг байдаг. Нөгөө талаас, анги дотор суралцагч бүрийн ойлголтын түвшин харилцан адилгүй байдаг тул багш бүх оюутанд нэгэн зэрэг, нэг түвшний тайлбар өгөхөөс өөр аргагүй нөхцөл түгээмэл тохиолддог. Үүний улмаас зарим оюутан агуулгыг бүрэн ойлголгүй дараагийн сэдэв рүү шилжих, эсвэл өөрт тохирсон давтлага, тайлбар авч чадалгүй хоцрох асуудал үүсдэг.

Уламжлалт \eng{e-learning} системүүд энэ асуудлыг тодорхой хэмжээнд шийдэхийг зорьдог боловч тэдгээрийн ихэнх нь контент хадгалах, файл түгээх, шалгалт авах, даалгавар өгөх зэрэг үндсэн үйлдлүүдэд төвлөрсөн байдаг. Өөрөөр хэлбэл, ийм системүүд нь сургалтын материалыг хэрэглэгчид хүргэх боломжтой боловч суралцагч тухайн \\ материалыг хэр зэрэг ойлгож байгаа, ямар хэсэг дээр хүндрэлтэй байгаа, ямар тайлбар хэрэгтэй байгааг нарийн түвшинд мэдэрч ажиллах боломж хязгаарлагдмал байдаг. Мөн оюутан өөрийн оруулсан эсвэл хэрэглэж буй материалаас шууд асуулт асууж, баталгаатай эх сурвалж дээр үндэслэсэн тайлбар авах боломж ихэнхдээ хангалтгүй байдаг. Иймээс зөвхөн агуулга түгээдэг системээс илүү, агуулгыг ойлгоход тусалдаг, тайлбарладаг, чиглүүлдэг, суралцагчийн ахицыг харгалздаг ухаалаг системийн хэрэгцээ өсөн нэмэгдэж байна.

Энэ хэрэгцээг хангах боломжит чиглэлийн нэг нь \eng{AI-powered learning assistant} юм. Ийм систем нь суралцагчийн асуултад хариулах, хичээлийн материалыг тайлбарлах, ойлгоход хүндрэлтэй хэсгийг энгийнээр задлах, цаашид давтах шаардлагатай сэдвийг санал болгох зэрэг үйлдлээр дамжин сургалтын үйл явцыг илүү идэвхтэй дэмжих боломжтой. Ялангуяа сүүлийн жилүүдэд \eng{Large Language Model (LLM)}-уудын хөгжил нь боловсролын орчинд асуулт-хариулт, тайлбар, нэгтгэл, зөвлөмж өгөх шинэ боломжуудыг нээж өгсөн. Гэвч ерөнхий \eng{LLM}-д суурилсан системүүдийн сул тал нь хариулт нь үргэлж хэрэглэгчийн бодит сургалтын материалд суурилдаггүй, зарим тохиолдолд баталгаагүй эсвэл хэт ерөнхий мэдээлэл өгөх эрсдэлтэй байдаг.

Ийм сул талыг бууруулахад \eng{Retrieval-Augmented Generation (RAG)} технологи онцгой ач холбогдолтой. \eng{RAG} нь системийг зөвхөн загварын дотоод мэдлэг дээр ажиллуулахгүй, харин хэрэглэгчийн оруулсан баримт бичиг, сургалтын материал, эх сурвалжаас холбогдох хэсгийг эхэлж олж авч, түүнд тулгуурлан хариулт үүсгэх боломж олгодог. Ингэснээр хариулт нь илүү эх сурвалжид суурилсан, илүү найдвартай, хэрэглэгчийн бодит хэрэгцээнд ойр болдог. Боловсролын орчинд энэ нь онцгой давуу талтай. Учир нь оюутан өөрийн судалж буй лекц, ном, гарын авлага дээр тулгуурлан асуулт асууж, яг тэр агуулгын хүрээнд тайлбар авах боломжтой болдог. Мөн ийм систем нь citation буюу эх сурвалжийн холбоос, эшлэлийг хамт үзүүлэх замаар хариултын баталгаатай байдлыг нэмэгдүүлж чадна.

Үүнээс гадна боловсролын салбарт хувь хүнд тохирсон сургалтын ач холбогдол улам бүр нэмэгдэж байна. Суралцагч бүрийн мэдлэгийн түвшин, ойлголтын хурд, сонирхол, сул болон давуу тал харилцан адилгүй байдаг тул ижил агуулгыг бүх хүнд ижил хэлбэрээр хүргэх нь үргэлж хамгийн оновчтой арга биш юм. Харин суралцагчийн аль сэдвийг сайн эзэмшсэн, аль сэдэв дээр илүү дэмжлэг хэрэгтэй байгааг илрүүлж, түүнд тохирсон тайлбар, зөвлөмж, давтлага өгөх нь сургалтын үр өгөөжийг нэмэгдүүлэх боломжтой. Энэ утгаараа \eng{personalized learning} буюу хувь хүнд тохирсон сургалт нь зөвхөн технологийн шинэ боломж бус, орчин үеийн боловсролын зайлшгүй хэрэгцээ болж байна.

Ийм үндэслэлүүдийн хүрээнд энэхүү дипломын ажил нь суралцагчийн оруулсан сургалтын материалд тулгуурлан ажиллах, асуултад эх сурвалжтай хариулах, мөн суралцагчийн мэдлэгийн ахицыг хянах боломжтой Суралцагчийн ухаалаг туслах системийн судалгаа, архитектур, хэрэгжүүлэлтийг авч үзэж байна. Тус систем нь \eng{RAG}, \eng{LLM}, \eng{Learner Model}, \eng{Mastery Tracking} зэрэг ойлголтыг нэгтгэн, зөвхөн асуултад хариулах бус, суралцах үйл явцыг илүү дэмжих зорилготой ухаалаг сургалтын туслах орчин бүрдүүлэхэд чиглэж байна. Иймээс энэ ажил нь боловсролын технологийн практик хэрэгцээ, хиймэл оюунд суурилсан шийдлийн боломж, хувь хүнд тохирсон сургалтын ач холбогдлыг нэгтгэсэн судалгааны ач холбогдолтой гэж үзэж байна.

\end{abstract}